\title{Parameter-Efficient Orthogonal Finetuning via Factorization}

\begin{abstract}
With the widespread use of large foundation models, the financial and computational burdens of training them from the ground up have become prohibitive. Consequently, methods that facilitate efficient adaptation of these sophisticated models to new tasks have grown in importance. This work investigates a systematic finetuning approach, Orthogonal Finetuning (OFT), for adapting models to downstream tasks. Although OFT exhibits robust generalization across applications, it incurs substantial parameter overhead due to the inherently high dimensionality of orthogonal matrices. To mitigate this, we approach OFT through the lens of information transmission, outlining core criteria for enhanced parameter-efficiency. Drawing inspiration from the Cooley-Tukey fast Fourier transform, known for its communication efficiency, we introduce a compact orthogonal parameterization via butterfly structures. Incorporating this innovation into OFT, we present a novel method—Orthogonal Butterfly (BOFT)—that dramatically reduces parameter usage for finetuning. BOFT generalizes OFT by treating it as a special scenario within a broader orthogonal adaptation framework. Finally, we comprehensively evaluate BOFT in adapting large-scale vision transformers, expansive language models, and text-to-image diffusion architectures, demonstrating its effectiveness on a diverse array of vision and language downstream tasks.
\end{abstract}

\section{Introduction}

Modern models such as ChatGPT~\cite{brown2020language,chen2021evaluating} and Stable Diffusion~\cite{rombach2022high} showcase the extraordinary generalization capacities of large foundation models. However, as model performance continues to soar, so too does the parameter count—GPT-3, for example, contains approximately 175B parameters—which makes full model training increasingly infeasible for most research groups. Sustained progress, therefore, hinges on the ability to adapt pre-existing foundation models to specific tasks without the prohibitive expense of full retraining. Efficient adaptation is commonly achieved via three principal strategies: \emph{model finetuning}~\cite{hulora2022,zaken2022bitfit,guo2020parameter,raffel2020exploring,qiu2023controlling,zhang2022adaptive,chavan2023one}, where a portion of the model's parameters are selectively updated without modifying the original architecture; \emph{adapter tuning}~\cite{rebuffi2017learning,houlsby2019parameter,pfeiffer2020adapterfusion,lin2020exploring,he2021towards}, which augments the model with additional trainable components; and \emph{prompt tuning}~\cite{li2021prefix,lester2021power}, which prepends learnable tokens to the model's input. Of these, model finetuning stands out for its conceptual simplicity and, crucially, its ability to introduce no extra inference delay.

The central concept of model finetuning revolves around maintaining similarity between the pretrained and finetuned models, as quantified by particular metrics, in order to retain the information acquired during pretraining. In practice, many contemporary finetuning approaches implement a low learning rate, which serves as an informal means to keep the difference between the original and adapted models minimal. Considering the inherent structure present in weight matrices, some more systematic strategies aim to safeguard the relational structure among weights, such as preserving the angular relationships between neurons. Building upon this idea, Orthogonal Finetuning~(OFT)~\cite{qiu2023controlling} introduces an innovative finetuning paradigm in which all neurons within a layer are subjected to the same orthogonal transformation, thereby ensuring that their mutual angles remain unchanged throughout the adaptation process. OFT has shown notable improvements in generalization and optimization when applied to text-to-image diffusion model finetuning~\cite{qiu2023controlling}; however, the orthogonal matrices involved are typically high-dimensional, which leads to a substantial number of trainable parameters. To enhance efficiency, OFT leverages a block-diagonal matrix structure to decrease parameter count. This solution, nevertheless, introduces trade-offs—namely, it fixes the sparsity pattern and imposes the orthogonal transformation independently within each block. While this block-diagonal configuration economizes on parameters, it can inadvertently inject restrictive inductive biases (for instance, diminishing model expressiveness and hindering approximation of standard linear transformations). Furthermore, determining an optimal sparsity configuration for such matrices remains an open question.

The central challenge lies in constructing a dense orthogonal matrix while maintaining parameter efficiency. At first glance, this seems prohibitive, as a dense orthogonal matrix in $d$ dimensions typically demands $\mathcal{O}(d^2)$ parameters. Our strategy deviates from convention by assembling a dense orthogonal matrix through the composition of several sparse, factorized matrices. The guiding principle is that computational resources can be exchanged for memory efficiency—while we must perform repeated multiplications of sparse matrices during each training step, the total count of trainable parameters is substantially reduced. We contextualize the matrix factorization process as an instance of information transmission, where the synthesis of a dense orthogonal matrix from sparse factors is analogous to the transmission of information across a grid-shaped graph structure. This perspective paves the way for numerous sparse matrix factorization schemes that strategically cap the number of learnable parameters yet retain the expressive power necessary for generating dense matrices.

In seeking further parameter reduction within this information transmission paradigm, we draw inspiration from the butterfly patterns in the Cooley-Tukey fast Fourier transform algorithm~\cite{cooley1965algorithm}, noted for their efficient information propagation among nodes~\cite{klasing1994broadcasting}. The butterfly graph pattern intrinsic to the Cooley-Tukey method informs an approach to sparse matrix decomposition, referred to as butterfly factorization. Through this, a $d$-dimensional dense matrix can be constructed via a sequence of $\mathcal{O}(\log d)$ sparse matrices, where each has $\mathcal{O}(d)$ non-zero elements. If every sparse factor remains orthogonal, this leads to a highly efficient orthogonal matrix parameterization, utilizing only $\mathcal{O}(d\log d)$ parameters—a drastic improvement over the classical approach. Building on this efficient structure, we introduce Orthogonal Butterfly~(BOFT), a new method for parameter-efficient finetuning. BOFT generalizes OFT as a special case and establishes a comprehensive orthogonal finetuning framework. A key point of commonality between block-diagonal and butterfly structures is their inherent sparsity; both embed sparse patterns within orthogonal matrices, thus curtailing the number of parameters requiring training. In comparing our method to the low-rank approach in LoRA~\cite{hulora2022}, it is worth noting that both low-rank and sparse matrices constitute principal classes of structured matrices~\cite{candes2011robust} leveraged for parameter efficiency.

In contrast to the block-diagonal approach employed by OFT, which attempts to balance expressivity and regularity, BOFT leverages the butterfly structure to achieve a more continuous spectrum between matrices drawn from the full orthogonal group (maximizing expressivity) and identity matrices (maximizing regularity). This construction allows us to delineate a narrower, more manageable hypothesis space within the orthogonal group, tailored to downstream learning objectives. Given the prevalence of the butterfly structure in numerous efficient linear operations—such as the discrete Fourier and discrete cosine transforms—we posit that the incorporation of such structured parameterizations not only increases parameter efficiency but also embeds a useful inductive bias that enhances model generalization and curbs overfitting. Our rationale is informed by the unique ability of the butterfly structure to represent many classical linear transforms—an expressive capacity unattainable by the block-diagonal structures used in OFT. Our main contributions are as follows:

\begin{itemize}[leftmargin=*,nosep]
    \item We approach the challenge of parameter efficiency in orthogonal finetuning via a new information transmission paradigm, recasting the optimization of a compact dense orthogonal matrix as the problem of transmitting information through a grid-based graph.
    \item Drawing inspiration from butterfly factorizations, such as those used in the Cooley-Tukey algorithm, we introduce Orthogonal Butterfly, a new orthogonal finetuning technique developed within our proposed information transmission framework.
    \item We provide several theoretical arguments explaining how BOFT can dramatically cut down the count of learnable parameters without sacrificing expressive power or training stability. Additionally, the matrix factorization inherent to BOFT facilitates a novel perspective on weight interpolation (see Figure~\ref{fig:boft_interpolation}).
    \item We are the first to deploy orthogonal finetuning across a spectrum of tasks beyond controllable text-to-image generation~\cite{qiu2023controlling}, demonstrating its broad applicability as a general finetuning strategy. Accordingly, we evaluate BOFT in adaptation scenarios spanning both computer vision and natural language processing. Empirically, BOFT substantially exceeds the performance of existing parameter-efficient approaches, affirming its superior efficiency and generalization capability.
\end{itemize}

\section{Related Work}

\textbf{Parameter-efficient finetuning (PEFT)}. As foundation models (\emph{e.g.}, \cite{brown2020language,radford2021learning,rombach2022high,kirillov2023segment}) continue to grow in both size and capability, efficiently adapting these models for specific downstream tasks has attracted significant attention~\cite{treviso2023efficient,ding2023parameter,lialin2023scaling}. Numerous PEFT strategies have been developed~\cite{houlsby2019parameter,aghajanyan2020intrinsic,hulora2022,edalati2022krona,wang2022adamix,gheini2021cross,zaken2022bitfit,guo2020parameter,sung2021training,ansell2022composable,lester2021power,li2021prefix,vu2022spot,he2021towards,mao2021unipelt,karimi2021compacter,liu2022few,sung2022lst,chen2022parameter,jia2022visual,chen2022adaptformer,zhang2022neural,jie2023fact,lian2022scaling,luo2023towards}, but approaches based on reparameterization~\cite{aghajanyan2020intrinsic,hulora2022,edalati2022krona} are particularly pertinent to our approach. LoRA~\cite{hulora2022} introduces a low-rank decomposition to the pretrained weight matrices, yielding notable success on NLP tasks by augmenting weights with the product of two low-rank matrices. Since LoRA assigns a fixed rank to these matrices, subsequent works~\cite{zhang2022adaptive,valipour2022dylora,zhang2023increlora} have investigated adaptive ranking, allowing variable allocation across layers to optimize parameter usage. The straightforwardness of low-rank weight reparameterization has led to its widespread adoption~\cite{dettmers2023qlora,zi2023delta,chavan2023one}. Drawing from research on hyperspherical energy and its implications for generalization~\cite{liu2018learning,liu2021orthogonal}, \cite{qiu2023controlling} introduce orthogonal finetuning, offering an alternative paradigm for refining text-to-image diffusion models. In this scheme, OFT learns an orthogonal transformation within a layer, promoting better generalization and more reliable training dynamics compared to LoRA. However, this advantage is tempered by a generally higher number of trainable parameters in OFT relative to LoRA. Enhancing OFT’s parameter efficiency is, therefore, an important objective. Furthermore, it remains to be determined if OFT is suitable for more general adaptation tasks beyond the scope of text-to-image diffusion~\cite{qiu2023controlling}. BOFT addresses this challenge by leveraging butterfly factorization to compress OFT, substantially increasing its parameter efficiency and enabling a broader demonstration of orthogonal finetuning in diverse adaptation contexts.

\textbf{Butterfly structures}. The radix-2 Cooley-Tukey procedure decomposes an $N$-point discrete Fourier transform into a sequence of two $\frac{N}{2}$-point transforms, resulting in a distinctive butterfly pattern that may be formulated as a product of several sparse matrices, collectively referred to as a butterfly matrix. Butterfly matrices have been employed to parameterize orthogonal matrices, mitigating the need for pivoting in Gaussian elimination and enhancing computational efficiency~\cite{parker1995random}. Their use has also extended to improved stability in recurrent neural network training~\cite{jing2017tunable} and in kernel approximation tasks~\cite{munkhoeva2018quadrature}. Works such as~\cite{dao2019learning,dao2019kaleidoscope} exploit butterfly parametrizations to learn efficient linear transforms, while~\cite{chen2022pixelated,dao2022monarch} employ them to facilitate scalable neural network training. Applications of butterfly structures also encompass fast algorithms for matrix-vector multiplication~\cite{michielssen1996multilevel,o2010algorithm}, compact matrix representation~\cite{li2015butterfly}, and efficient network information dissemination~\cite{klasing1994broadcasting,li2003linear,rajasingh2016transmission}. Distinct from these previous investigations, our contribution is to integrate butterfly structures with OFT, substantially boosting parameter efficiency.

\section{An Information Transmission View on Orthogonal Finetuning}

We begin by reviewing the foundational concepts of OFT. In order to adapt a pretrained weight matrix, OFT expresses the updated weight matrix as the product of a trainable orthogonal matrix with the fixed pretrained weight matrix. In contrast to LoRA—which refines the weights via an additive, low-rank matrix—OFT employs a multiplicative update utilizing an orthogonal matrix. Both techniques pursue parameter efficiency: LoRA through enforcing a low-rank structure, and the initial formulation of OFT through a block-diagonal orthogonal matrix~\cite{qiu2023controlling}. Notably, reducing the block size of the diagonal blocks further enhances the parameter efficiency in OFT. Figure~\ref{fig:comp_lora} provides a visual comparison of these methodologies. 

The primary rationale for using an orthogonal transformation to refine the weight matrix is to maintain the angular relationships between neurons~\cite{liu2018learning,liu2021orthogonal,qiu2023controlling}, thus retaining much of the semantic knowledge acquired during pretraining. Specifically, OFT learns an orthogonal matrix $\bm{R} \in \mathbb{R}^{d \times d}$ associated with a pretrained linear transformation $\bm{W}^0 \in \mathbb{R}^{d \times n}$, thereby altering the layer computation from $\bm{z}=(\bm{W}^0)^\top\bm{x}$ to $\bm{z}=(\bm{R}\bm{W}^0)^\top\bm{x}$, with $\bm{x} \in \mathbb{R}^d$ and $\bm{z} \in \mathbb{R}^n$ representing the input and output, respectively. To guarantee zero initialization, OFT sets $\bm{R}$ to the identity matrix at the outset. 

Maintaining the orthogonality of $\bm{R}$ during finetuning is accomplished by adopting the Cayley parameterization, following the approach of~\cite{qiu2023controlling,liu2021orthogonal}: $\bm{R} = (\bm{I} + \bm{Q})(\bm{I} - \bm{Q})^{-1}$, where $\bm{Q}$ is skew-symmetric, that is, $\bm{Q} = -\bm{Q}^\top$. For increased parameter efficiency, a block-diagonal arrangement is used, where $\bm{R}$ is expressed as $\text{diag}(\bm{R}_1, \bm{R}_2, \ldots, \bm{R}_r)$ with each $\bm{R}_i \in \mathcal{R}^{b \times b}$ being a small orthogonal matrix and $br = d$. This design introduces efficiency at the expense of the assumption that a neuron's dimension (specifically, a column in $\bm{W}^0$) is segmented into $r$ groups, with distinct orthogonal matrices transforming each group. While the block-diagonal approach yields favorable empirical performance, partitioning neuron dimensions by their indices is not semantically meaningful, highlighting the appeal of dense orthogonal matrices. This raises an important question: \emph{Is it possible to construct a dense orthogonal matrix that retains parameter efficiency?}

To tackle this issue, we suggest expressing a dense orthogonal matrix as a product of several sparse orthogonal matrices. In order to establish a cohesive and accessible framework for analyzing the sparsity structures involved in orthogonal matrix decompositions, we recast the task of constructing a dense orthogonal matrix in OFT as analogous to an information transmission process. In detail, forming a dense matrix $\bm{R}\in\mathbb{R}^{d\times d}$ by multiplying $m$ square matrices, such that $\thickmuskip=2mu \medmuskip=2mu \bm{R}=\bm{B}_m\bm{B}_{m-1}\cdots\bm{B}_1$, can be understood as conveying information through a network represented by a $d\times (m+1)$ node grid (see Figure~\ref{fig:info_trans}). This information transmission analogy arises from the idea that a dense $d$-dimensional square matrix encapsulates complete interactions between two sets of $d$ nodes. For $\bm{R}$, the presence of a zero in position $R_{ij}$ signifies the absence of information flow from node $j$ to node $i$, whereas a nonzero entry permits the transfer of information. Accordingly, decomposing the dense matrix $\bm{R}$ into a sequence of matrices $\bm{B}_m\bm{B}_{m-1}\cdots\bm{B}_1$ equates to a staged transmission of information, dictated by the edge patterns of each graph induced by $\bm{B}_i$. The flow commences following the structure of $\bm{B}_1$ and terminates in accordance with $\bm{B}_n$. For illustration in Figure~\ref{fig:info_trans}, we examine the decomposition $\bm{R}=\bm{B}_5\bm{B}_4\bm{B}_3\bm{B}_2\bm{B}_1$ and depict both its sparsity configurations and the associated induced graph. This graph is constructed by expanding the sequence of matrix products. Within this induced graph, each $\bm{B}_i$ represents a connectivity matrix linking nodes from the $i$-th layer to the $(i+1)$-th layer. Specifically, element $(j_1,j_2)$ in $\bm{B}_i$ determines the existence of a directed edge from the $j_2$-th node in the $i$-th layer to the $j_1$-th node in the $(i+1)$-th layer; a zero value denotes no such connection. To guarantee that $\bm{B}_5\bm{B}_4\bm{B}_3\bm{B}_2\bm{B}_1$ yields a dense matrix, every node at the first layer must possess the ability to transmit to all nodes at the sixth layer. If we restrict attention to $\bm{R}=\bm{B}_4\bm{B}_3\bm{B}_2\bm{B}_1$, involving only the initial and fifth layer nodes, it can be observed that node $1$ fails to reach node $3$, indicating a zero at position $(3,1)$ in the combined sparsity pattern. This correspondence between the induced graphs and their underlying sparsity structures similarly applies to shorter products, such as $\bm{R}=\bm{B}_3\bm{B}_2\bm{B}_1$ and $\bm{R}=\bm{B}_2\bm{B}_1$. More generally, for any dense matrix $\bm{R}\in\mathbb{R}^{d\times d}$, the sequence of $m$ factor matrices $\bm{B}_m,\cdots,\bm{B}_1$ must define a collection of directed edges across a $d\times (m+1)$ grid—edges can only link adjacent layers (columns)—so that information originating from any node in the first layer can reach every node in the $(m+1)$-th layer.

The information transmission perspective offers valuable insights into the sparsity configuration found in the factorization matrices within OFT. As depicted in Figure~\ref{fig:blk_diag}, the original OFT’s $\bm{R}$ matrix exhibits a block-diagonal arrangement. While this architecture lowers the overall count of trainable parameters, it is inherently incapable of representing $\bm{R}$ as a fully dense matrix. Our objective is to construct a dense orthogonal matrix from $m$ constituent sparse orthogonal matrices, striving to minimize the effective number of trainable parameters required for this composition. According to the information transmission viewpoint, two primary principles must be satisfied: \emph{(i)} \textbf{dense connectivity}—each node from the input layer must be reachable by some path to every node in the output layer; and \emph{(ii)} \textbf{minimum free edges}—the overall number of edges should be minimized while upholding the orthogonality constraint. It is crucial to observe that orthogonality enforces intricate conditions on the connections linking adjacent levels. To illustrate, for any $\bm{B}_i$ to achieve full rank—a prerequisite for orthogonality—there must exist $d$ edges establishing a bijection between all nodes in the $i$-th layer and all nodes in the $(i=1)$-th layer. Consequently, the edge count between consecutive layers cannot fall below $d$ (see the example with 4 edges in Figure~\ref{fig:info_trans}). Since these $d$ edges serve to maintain orthogonality, and are not trainable (as for a $d\times d$ orthogonal matrix with $d$ non-zero entries, the only permissible values are $\pm 1$), they are excluded from the tally of trainable edges. The structural restrictions imposed by orthogonality thus reduce the number of required parameters compared to unconstrained matrices, introducing additional interdependence among edge placements. For instance, in a $2\times 2$ orthogonal matrix, if a zero occupies the $(1,1)$ entry, orthogonality dictates a corresponding zero in the $(2,2)$ entry, signifying that removing one edge might necessitate the removal of another. Hence, for a given admissible connection scheme, the orthogonality condition can necessitate either the inclusion or exclusion of certain edges. Through examination of the non-zero entries within the synthesized orthogonal matrix, the information transmission framework proves particularly pertinent to OFT, since the chief concern lies with the pattern of non-zero trainable components in $\bm{R}$ rather than their exact magnitudes.

A straightforward dense linkage between two layers necessitates $\mathcal{O}(d^2)$ connections (that is, a full dense orthogonal matrix), amounting to $d^2-d$ total connections (for $d=4$, this results in 12 connections). Figure~\ref{fig:info_trans} displays a possible matrix factorization that requires only $10$ connections, which is in fact fewer than what is required for a dense orthogonal matrix. This perspective allows for analysis of OFT's parameter efficiency from a graph-theoretical standpoint, and offers a systematic approach for generating viable factorizations. Our methodology is inspired by the butterfly graphs~\cite{cooley1965algorithm}, a structure arising in the Cooley-Tukey algorithm that provides efficient dense connectivity between $d$ input and $d$ output nodes using merely $\mathcal{O}(d\log d)$ links. For example, in Figure~\ref{fig:info_trans}, the configuration requires $10$ connections to establish full connectivity, whereas the butterfly network achieves the same with just $8$ connections. In the following, we describe how employing the butterfly structure enhances parameter efficiency.

\section{Orthogonal Parameterization by Butterfly Factorization}

The butterfly network, foundational in the Cooley-Tukey fast Fourier transform algorithm, is designed to facilitate efficient computation. In the context of the Fourier transform, localized alterations in the frequency domain translate to global modifications in the spatial domain, which parallels our problem of information transfer—where any node at the initial level is capable of passing information to every node at the terminal level. Moreover, the butterfly network topology~\cite{solihin2015fundamentals,li2011linear} has found widespread use in computer networks to optimize information flow. Consider $k\geq 2$ as a power of $2$, and introduce the \emph{butterfly factor} $\bm{B}^F(k)$ as

\begin{equation}\label{eq:but_factor}
\begin{aligned}
    \bm{B}^F(k)=\begin{bmatrix}
    \text{diag}(\bm{d}_1) & \text{diag}(\bm{d}_2) \\
    \text{diag}(\bm{d}_3) & \text{diag}(\bm{d}_4)
    \end{bmatrix}\in\mathbb{R}^{k\times k},
\end{aligned}
\end{equation}

with $\bm{d}_i\in\mathbb{R}^{\frac{k}{2}}$ for all $i$, representing appropriate vectors. If we let $d=2^N$, the $d$-dimensional \emph{butterfly matrix} $\bm{B}(d)\in\mathbb{R}^{d\times d}$ is then defined recursively as follows:

\begin{equation}
\begin{aligned}
    \!\!\bm{B}(d)=\tilde{\bm{B}}(d,d)\!\cdot\!\begin{bmatrix}
    \bm{B}_1(\frac{d}{2})\!\!\!\!\! & \bm{0} \\
    \bm{0}\!\!\!\!\! &\bm{B}_2(\frac{d}{2})
    \end{bmatrix}= \tilde{\bm{B}}(d,d)\tilde{\bm{B}}(d,\frac{d}{2})\cdots\tilde{\bm{B}}(d,2),
\end{aligned}
\end{equation}

where $\bm{B}_1(\frac{d}{2})$ and $\bm{B}_2(\frac{d}{2})$ denote two butterfly matrices each of dimension $\frac{d}{2}$. We introduce the notion of a \emph{butterfly component} by setting $\thickmuskip=2mu \medmuskip=2mu \tilde{\bm{B}}(d,k)=\text{diag}(\bm{B}^F_1(k),\cdots,\bm{B}^F_{d/k}(k))$, representing a block-diagonal matrix of dimensions $d\times d$ where each block has size $k$. Here, the $\bm{B}^F_i(k)$ for all $i$ correspond to the butterfly factors described in Equation~\ref{eq:but_factor}. With this setup, we can parameterize orthogonal matrices using butterfly matrices. Specifically, ensuring that each multiplicative term $\tilde{\bm{B}}(d,k)$ within the butterfly matrix $\bm{B}(d)$ is orthogonal suffices. Focusing on the block-diagonal matrix $\tilde{\bm{B}}(d,2)$ with blocks of size $2$, orthogonality is readily enforced by applying the Cayley transform (i.e., $2$-dimensional rotations) to each block, as established in \cite{liu2021orthogonal,qiu2023controlling}. The non-zero structure present in any butterfly component is simply a permutation of the pattern in $\tilde{\bm{B}}(d,2)$, making it straightforward to extend the orthogonal parameterization to all butterfly components. As a result, this method yields an efficient way to construct orthogonal matrices using numerous $2\times2$ orthogonal blocks. Extending this construction according to \cite{chen2022pixelated}, we define a block butterfly component $\tilde{\bm{B}}^b(d,k)$, in which each element of $\bm{d}_i$ for all $i$ becomes a $b\times b$ matrix. To ensure that $\tilde{\bm{B}}^b(d,2)$ is orthogonal, every $2b\times2b$ block must itself be orthogonal. The non-zero structure for any $\tilde{\bm{B}}^b(d,k)$, with $k>2$, results from a block-wise permutation of the non-zero structure in $\tilde{\bm{B}}^b(d,2)$, and can therefore also be parameterized as an orthogonal matrix using similar methods. Putting these elements together, the forward computation in BOFT is

\begin{equation}
    \begin{aligned}\nonumber
    \bm{z}=\big(\bm{R}(m,b)\cdot\bm{W}^0\big)^\top\bm{x},~~~~\text{s.t.}~\bigg{\{}\bm{R}(m,b)=\prod_{i=1}^m \tilde{\bm{B}}^b_{(d,i)}~~\&~~\underbrace{\big(\tilde{\bm{B}}^b_{(d,j)}\big)^\top\tilde{\bm{B}}^b_{(d,j)}=\tilde{\bm{B}}^b_{(d,j)}\big(\tilde{\bm{B}}^b_{(d,j)}\big)^\top\!=\bm{I}_d}_{\forall j\in [1, m]}\bigg{\}},
    \end{aligned}
\end{equation}

In this formulation, for notational convenience, we let $\tilde{\bm{B}}^b(d,2^{m - i+1})$ be represented as $\tilde{\bm{B}}^b_{(d,i)}$, with $\bm{I}_d$ signifying the $d$-dimensional identity matrix. The orthogonal transformation $\bm{R}(m,b)\in\mathbb{R}^{d\times d}$ is constructed by a sequence of orthogonal butterfly matrices. Throughout the rest of the paper, we refer to BOFT parametrized with $\bm{R}(m,\frac{b}{2})$ as BOFT($m$,$b$), for $b\geq 2$. In particular, when $m=1$, BOFT($1$,$b$) reduces to the block-diagonal OFT~\cite{qiu2023controlling}, using a block size of $b$. Further, BOFT($1$,$d$) corresponds to the standard OFT~\cite{qiu2023controlling} with a full, unconstrained orthogonal matrix. Setting $m=\log \frac{2d}{b}$ yields BOFT($\log \frac{2d}{b}$,$b$), which utilizes the block butterfly matrix $\bm{B}^{\frac{b}{2}}(d)$ as $\bm{R}$, resulting in a dense orthogonal matrix. Generally, the number of trainable parameters in BOFT($m$,$b$) amounts to $\frac{1}{2}(b-1)dm$ for a linear layer of shape $d\times n$. Specifically, if $m=\log d$ and $b=2$ (i.e., using butterfly matrices), BOFT requires $\mathcal{O}(d\log d)$ parameters. In comparison, conventional OFT with a full dense orthogonal matrix needs $\mathcal{O}(d^2)$ parameters, and block-diagonal OFT with $r$ blocks consumes $\mathcal{O}(bd)$. As a result, to realize a dense orthogonal matrix, original OFT must take $b=d$ as the block size, while BOFT can achieve this density using arbitrary $b$.

\textbf{Identity Initialization in BOFT}. Typical finetuning strategies aim to ensure minimal deviation from the initial pretrained model parameters. For instance, LoRA employs zero-initialization for the low-rank update matrices. Similarly, in BOFT, each butterfly factor is initialized as the identity matrix (that is, the underlying skew-symmetric matrix within the Cayley parameterization starts as zero).

\textbf{Multiplicative Dropout within BOFT}. In analogy to LoRA~\cite{hulora2022}, which uses Dropout for its low-rank components to regularize training, we observe that classical Dropout~\cite{srivastava2014dropout} does not fit BOFT due to its multiplicative update scheme. To address this, we design a multiplicative Dropout mechanism specific to BOFT. Since $\bm{R}(m,b)$ is formed from $m$ butterfly matrices, each of which can be expressed as block-diagonal orthogonal matrices with block size $2b\times 2b$, our multiplicative Dropout randomly selects $p_1$ fraction of butterfly matrices and $p_2$ proportion of blocks within each, and substitutes the selected blocks with identity matrices.

\section{Intriguing Insights and Discussions}

\textbf{Expressivity of BOFT}. The butterfly architecture, augmented by permutation operations, is known to exactly reconstruct several standard fast linear transforms~\cite{dao2019learning,dao2019kaleidoscope} (\emph{e.g.}, the fast Fourier transform and the Hadamard transform). Nonetheless, the capacity of the orthogonal butterfly matrix to approximate arbitrary orthogonal matrices has yet to be fully characterized. To shed light on this, we perform experiments attempting to approximate a randomly sampled dense orthogonal matrix~\cite{anderson1987generation} of dimension $512\times 512$. As shown in Figure~\ref{fig:para_eff}, all outcomes are averaged across 10 independent random initializations. The vertical axis measures approximation error, while the horizontal axis reflects the number of effective learnable parameters. Each color traces the results for BOFT with a fixed block size, and the leftmost value corresponds to the approximation error for BOFT($1$,$b$) (\emph{i.e.}, the conventional block-diagonal OFT of block size $b$). Broadly speaking, BOFT consistently achieves greater parameter efficiency compared to OFT. For instance, BOFT($9$,$2$) delivers higher expressivity than BOFT($1$,$16$), with a markedly reduced parameter count. Parameter efficiency further improves as the block size $b$ decreases and the butterfly depth $m$ increases. Illustratively, BOFT($6$,$4$) attains a similar approximation error as BOFT($2$,$16$) but utilizes significantly fewer parameters. Overall, the butterfly matrix corresponds to a richer and more structured subset of the orthogonal group than does the block-diagonal arrangement, rendering BOFT provably more expressive than OFT for a fixed block size.

\begin{theorem}[Expressivity of BOFT]\label{thm:exp_boft}
    BOFT surpasses OFT in expressivity when the block size is held constant. Any orthogonal matrix of dimension $d$ can be approximated by a product of butterfly matrices of the form $\bm{B}_{d-1,1}(d)\bm{B}^\top_{d-1,2}(d)\cdots\bm{B}_{1,1}(d)\bm{B}^\top_{1,2}(d)$, where each $\bm{B}_{i,j}(d)$ is a butterfly matrix.
\end{theorem}

According to Theorem~\ref{thm:exp_boft}, there is a straightforward route to generalize BOFT: The general form of the target orthogonal matrix becomes $\thickmuskip=2mu \medmuskip=2mu \bm{R}^G(m_1,b_1,m_2,b_2,l)=\bm{R}_{l,1}(m_1,b_1)\bm{R}_{l,2}^T(m_2,b_2)\cdots\bm{R}_{1,1}(m_1,b_1)\bm{R}_{1,2}^T(m_2,b_2)$, with $\bm{R}_{i,j}^T(m,b)$ indicating the orthogonal matrices within BOFT. In the specific case where $m_1=m_2=\log d$, $b_1=b_2=2$, and $l=d-1$, this form, $\bm{R}^G(m_1,b_1,m_2,b_2,l)$, spans the entire orthogonal group. Such composite structures are often referred to as kaleidoscope hierarchies~\cite{dao2019kaleidoscope}. Nevertheless, it is important to observe that greater expressivity does not inherently translate into better practical outcomes during finetuning: universal expressivity, as seen in full finetuning, frequently fails to yield optimal results. Striking a balance between expressivity and inductive bias is central for successful generalization in model finetuning. BOFT leverages structural priors to expand the range of trainable parameters, facilitating a superior trade-off between flexibility and regularization.

\textbf{Spectral properties}. Compared to LoRA, orthogonal finetuning generally offers superior preservation of spectral properties, as it exactly maintains the spectral norm of the original pretrained weight matrix $\bm{W}^0$. This can be formally understood via singular value decomposition: given $\bm{W}^0=\bm{U}\bm{\Sigma}\bm{V}^\top$, where $\bm{U}$ and $\bm{V}$ are orthogonal matrices and $\bm{\Sigma}$ contains the singular values along its diagonal, both OFT and BOFT apply an orthogonal transformation $\bm{R}$ from the left, resulting in updated weights $\bm{R}\bm{U}\bm{\Sigma}\bm{V}^\top$. This operation leaves the largest singular value unchanged, thereby preserving the spectral norm of $\bm{W}^0$. Such invariance has been shown to be advantageous for maintaining training stability and improving generalization performance~\cite{miyato2018spectral,yoshida2017spectral}. Additional mathematical implications are explored in Appendix~\ref{app:sn_property}.

\textbf{Orthogonal finetuning as learning bilinear similarity}. The BOFT method can be interpreted as learning a bilinear similarity of the form $\bm{w}^0_i\bm{R}\bm{x}$, where $\bm{w}^0_i$ denotes the $i$-th neuron (specifically, a column vector) from $\bm{W}^0$. In this framework, BOFT effectively learns a bilinear similarity transformation matrix $\bm{R}$ under the restriction that $\bm{R}$ remains orthogonal. This constraint links the approach conceptually to distance metric learning~\cite{xing2002distance} and the theory of bilinear forms~\cite{roman2005advanced}.

\textbf{Inductive bias and generalization in BOFT}. In BOFT, $\bm{R}(m,b)$ typically spans a structured subset of the orthogonal group, which imposes limitations on the class of hypotheses considered and thereby introduces an inductive bias. The specific form of inductive bias induced by the butterfly factorization is hypothesized to enhance generalization, given its shared architectural patterns with classical linear transforms~\cite{dao2019kaleidoscope}, such as the discrete Fourier, sine/cosine, and Hadamard transforms. Furthermore, the utilization of sparse matrix factorization in BOFT is likely to impart an additional implicit inductive bias~\cite{gunasekar2017implicit,li2018algorithmic,liu2019neural}.

\textbf{Comparison to butterfly-based sparse training}. Numerous studies~\cite{dao2019learning,dao2019kaleidoscope,chen2022pixelated,dao2022monarch} have examined sparse training utilizing the butterfly parameterization. These works often aim to directly reparameterize weight matrices with the butterfly structure, enabling the training of neural networks from initialization. In~\cite{dao2022monarch}, the approach is to fine-tune pretrained weights by first mapping them onto a variation of butterfly matrices, then optimizing these projected representations for subsequent tasks. In contrast, BOFT introduces a distinct fine-tuning methodology, in which the transformation of weights employs layer-shared weight matrices.

\input{rebuttal-results}

\section{Concluding Remarks and Limitations}

We introduce Orthogonal Butterfly, a broadly applicable and parameter-efficient finetuning framework for foundation models that leverages the butterfly architecture. The central idea for achieving enhanced parameter efficiency is to represent a dense orthogonal matrix as a product of several sparse orthogonal matrices. To facilitate the discovery of valid matrix decompositions, we develop a graphical framework for information transmission. Using this perspective, we demonstrate that the butterfly architecture is well-suited for our goal of sparse orthogonal matrix decomposition. Empirical results showcase BOFT’s strong performance in adapting large language models, large vision models, and text-to-image generation models. Our experimental findings further substantiate the effectiveness of BOFT as a universal approach for model finetuning.

However, BOFT is not without drawbacks. Specifically, because BOFT represents the final orthogonal transformation as a sequence of orthogonal matrices, the training phase incurs slightly higher runtime costs compared to OFT. Addressing the training efficiency of BOFT thus remains an unresolved issue. Fortunately, upon completing the finetuning process, the learned BOFT orthogonal matrices can be fused into the original model weights, thereby eliminating any extra inference cost. Additionally, it remains unclear whether the butterfly topology provides the optimal means for information transfer. Our framework for information transmission opens avenues for further exploration, taking cues from the field of computer networking, which extensively studies optimal network structures for information dissemination. We anticipate that alternative network architectures may enable even more efficient constructions of dense orthogonal matrices.

\end{document}